\section{CLI \texttt{sk}}\label{sec:cli}

Każdy krok łańcucha działa z~wiersza poleceń, w~pipeline (zasada projektowa 7). Rozdział
opisuje każde polecenie \repopath{sk} z~każdą opcją, odpowiadającym jej kluczem
\repopath{secondkey.yaml}, wartością domyślną i~zachowaniem, kody wyjścia, pełny plik
konfiguracji, demonstrację end-to-end na sample shop i~sposób czytania wyniku. Formaty plików,
które polecenia czytają i~piszą, opisuje sekcja~\ref{sec:formaty}; implementację ---
sekcja~\ref{sec:rdzen}.

\subsection{Wywołanie i~reguły wspólne}\label{sec:cli-wywolanie}

\repopath{sk} jest pakowany jako narzędzie .NET (\repopath{PackageId} \repopath{SecondKey.Cli},
polecenie \repopath{sk}), ale nie został jeszcze opublikowany:

\begin{lstlisting}
dotnet tool install --global SecondKey.Cli   # once published; until then:
dotnet run --project src/SecondKey.Cli -- --help
\end{lstlisting}

{\small
\begin{tabularx}{\linewidth}{@{}>{\raggedright\arraybackslash}p{3.4cm} >{\raggedright\arraybackslash}X >{\raggedright\arraybackslash}p{4cm}@{}}
\toprule
Polecenie & Krok & Stan w~tej wersji \\
\midrule
\texttt{sk capture} & C1 --- nagrywa ruch przez proxy przed systemem legacy & dostępne (S9) \\
\texttt{sk mine} & C2 --- proponuje klauzule kontraktu z~nagrań & planowane, faza 02
(kod 70) \\
\texttt{sk replay} & C6 --- odtwarza capture na obu stronach w~identycznych warunkach &
dostępne (S10–S11) \\
\texttt{sk compare} & C7 --- porównuje obie strony względem kontraktu do
\texttt{verdict.json} & dostępne (S12–S13) \\
\texttt{sk gate} & C5 --- podsumowuje SARIF bramki per reguła & dostępne (S14) \\
\texttt{sk mutate} & C9 --- wstrzykuje defekty, które kontrakt musi wykryć & planowane,
faza 02 (kod 70) \\
\texttt{sk evidence} & C10 --- składa pakiet dowodowy & dostępne (S14) \\
\texttt{sk validate <file>...} & sprawdza artefakty względem ich formatów & dostępne \\
\bottomrule
\end{tabularx}
}

\paragraph{Konfiguracja i~pierwszeństwo.} Opcja globalna \repopath{-{}-config <path>} wskazuje
każdemu poleceniu plik \repopath{secondkey.yaml}; domyślnie jest to \path{./secondkey.yaml},
jeśli istnieje. Plik podany jawnie, którego nie ma, jest błędem konfiguracji (kod 3). Opcje z~wiersza poleceń mają pierwszeństwo przed plikiem; wyjątkiem są dwie opcje listowe
\repopath{sk capture}, \repopath{-{}-session-key} i~\repopath{-{}-redact-header}, które
dopisują się do list z~pliku. Ścieżka podana w~opcji jest rozwiązywana względem bieżącego
katalogu, a~ścieżka z~pliku --- względem katalogu, w~którym leży plik. Polecenia poza
\repopath{sk validate} czytają konfigurację na początku, więc zepsuty plik jest zgłaszany
niezależnie od uruchomionego polecenia. Błąd składni wiersza poleceń kończy się komunikatem
„Run 'sk -{}-help' for usage.” i~kodem 64.

\zrodla{\path{letsgolegacy.secondkey/docs/cli.md},
\path{letsgolegacy.secondkey/src/SecondKey.Cli/SecondKeyCli.cs},
\path{letsgolegacy.secondkey/src/SecondKey.Cli/SecondKey.Cli.csproj},
\path{letsgolegacy.secondkey/src/SecondKey.Cli/Commands/StubCommand.cs},
\path{letsgolegacy.secondkey/src/SecondKey.Cli/Configuration/SecondKeyConfig.cs}}

\subsection{\texttt{sk capture}}\label{sec:cli-capture}

\begin{lstlisting}
sk capture --target http://127.0.0.1:5080 --listen http://127.0.0.1:8080 --session-key ASP.NET_SessionId
\end{lstlisting}

Uruchamia proxy nagrywające. Klienci łączą się z~adresem \repopath{-{}-listen} zamiast z~systemem legacy; każde żądanie trafia do \repopath{-{}-target} bez zmian --- poza tym, że
proxy prosi o~nieskompresowaną odpowiedź, żeby nagranie dało się czytać --- a~klient dostaje
odpowiedź systemu legacy bez zmian. \finding{Systemu legacy nie dotyka się: zmienia się tylko
adres, pod który łączą się jego klienci.}

Rysunek~\ref{fig:capture} pokazuje drogę jednej wymiany przez proxy.

\begin{figure}[htbp]
  \centering
  \includegraphics[width=\linewidth]{\dgm{a3-sekwencja-capture}}
  \caption{Capture: klient, proxy \texttt{sk capture}, system legacy i plik nagrania.}
  \label{fig:capture}
\end{figure}

{\small
\begin{xltabular}{\linewidth}{@{}>{\raggedright\arraybackslash}p{4.9cm} >{\raggedright\arraybackslash}p{2.7cm} >{\raggedright\arraybackslash}X@{}}
\toprule
Opcja / klucz w~\texttt{secondkey.yaml} & Domyślnie & Znaczenie \\
\midrule
\endfirsthead
\toprule
Opcja / klucz w~\texttt{secondkey.yaml} & Domyślnie & Znaczenie \\
\midrule
\endhead
\bottomrule
\endlastfoot
\texttt{-{}-target <url>}\newline \path{capture.target} & brak --- wymagane & system legacy za
proxy \\
\texttt{-{}-listen <url>}\newline \path{capture.listen} & \path{http://127.0.0.1:8080} & gdzie
nasłuchuje proxy \\
\texttt{-{}-out <file>}\newline \path{capture.out} & \path{.secondkey/traffic.skcap} & nagranie \\
\texttt{-{}-session-key <name>}\newline \path{capture.sessionKeys} & brak & ciasteczko albo
nagłówek, którego wartość identyfikuje sesję użytkownika; powtarzalna \\
\texttt{-{}-redact-header <name>}\newline \path{capture.redactHeaders} & --- & nagłówek, którego
wartość nigdy nie trafia do nagrania, oprócz \texttt{authorization},
\texttt{proxy-authorization}, \texttt{cookie}, \texttt{set-cookie} i~\texttt{x-api-key},
redagowanych zawsze; powtarzalna \\
\texttt{-{}-max-body-bytes <n>}\newline \path{capture.maxBodyBytes} & 1 MiB & treść dłuższa jest
obcinana i~oznaczana jako obcięta \\
\texttt{-{}-exit-after <n>}\newline (tylko opcja) & brak & zatrzymanie po nagraniu \texttt{n} wymian \\
\end{xltabular}
}

\paragraph{Sesje i~scenariusze.} Klucz sesji decyduje o~scenariuszach: \repopath{sk replay}
odtwarza wymiany jednej sesji jako jeden scenariusz, po kolei. Wymiana należy do sesji, którą
nazywa jej żądanie --- ciasteczkiem albo nagłówkiem o~tej nazwie --- a~przy pierwszej wizycie
do tej, którą jej odpowiedź ustawia jako ciasteczko. Wymiana bez żadnego z~nich jest sesją
samą dla siebie, więc bez klucza każda wymiana jest odtwarzana osobno. W~nagraniu sesja jest
zapisana jako skrót, nigdy jako wartość ciasteczka.

\paragraph{Redakcja i~limit treści.} Wartości redagowanych nagłówków --- w~żądaniu i~w
odpowiedzi --- są zastępowane przez \repopath{<redacted>} przed zapisem na dysk, a~nagłówek
nagrania wymienia listę redagowanych nagłówków i~limit treści. Klient zawsze dostaje całą
odpowiedź; limit dotyczy tylko kopii w~nagraniu.

\paragraph{Gdy system legacy nie odpowiada.} Żądanie, na które system legacy nie odpowiedział
--- odrzucone, zerwane, przekroczony czas --- dostaje własną odpowiedź proxy 502 albo 504 i~\textbf{nie} jest nagrywane, bo błąd proxy nie jest zachowaniem systemu legacy. Takiej wymiany
brakuje potem w~jej scenariuszu, więc skryptowane nagranie przed użyciem pliku sprawdza, czy
liczba pominiętych wynosi 0. Nawiązanie połączenia z~systemem legacy ma limit 15 sekund.

\paragraph{Komunikaty i~kody wyjścia.} Po starcie polecenie wypisuje
\texttt{sk capture: recording <adres> -> <target> into <plik>. Stop with Ctrl+C.}, a~po
zatrzymaniu podsumowanie:

\begin{lstlisting}
sk capture: recorded 290 exchange(s), skipped 0 the legacy system did not answer.
\end{lstlisting}

Kod 0 po czystym zatrzymaniu; 64, gdy nie podano celu („no target --- pass -{}-target or set
capture.target in secondkey.yaml”) albo adres nie jest bezwzględnym URL http(s); 3 przy
błędzie konfiguracji.

\zrodla{\path{letsgolegacy.secondkey/docs/cli.md},
\path{letsgolegacy.secondkey/src/SecondKey.Cli/Commands/CaptureCommand.cs},
\path{letsgolegacy.secondkey/src/SecondKey.Capture/RecordingProxy.cs},
\path{letsgolegacy.secondkey/src/SecondKey.Capture/CaptureOptions.cs},
\path{letsgolegacy.secondkey/src/SecondKey.Capture/SessionTracker.cs}}

\subsection{Zatrzymywanie \texttt{sk capture}}\label{sec:cli-stop}

Ctrl+C albo SIGTERM zatrzymuje proxy czysto: nie przyjmuje ono nowych żądań, pozwala
dokończyć te w~toku --- najwyżej przez 30 sekund, potem je przerywa --- zamyka plik, wypisuje
podsumowanie i~kończy się kodem 0. Każde polecenie zatrzymane Ctrl+C albo SIGTERM ma minutę na
czyste zakończenie; domyślne dwie sekundy System.CommandLine byłyby za krótkie dla
\repopath{sk capture}, które czeka na wymiany w~toku do wolnego systemu legacy. Polecenie,
które nie zakończy się w~tym czasie, jest kończone kodem samego sygnału, 130 albo 143, a~to,
co zapisywało, może być niekompletne.

Każda wymiana jest zapisywana i~opróżniana z~bufora, gdy się zakończy, więc zabity capture
zachowuje to, co nagrał --- ale zabicie w~trakcie zapisu zostawia urwaną ostatnią linię,
którą \repopath{sk validate} odrzuca. Tam, gdzie pipeline może proces jedynie zabić (na
przykład runner Windows zatrzymujący zadanie w~tle), należy podać \repopath{-{}-exit-after} z~liczbą wymian, które wysyła skrypt, żeby capture zatrzymał się sam. Test z~sygnałem SIGTERM
sprawdza, że wymiana w~toku zostaje dokończona i~zapisana; na Windows jest pomijany.

\zrodla{\path{letsgolegacy.secondkey/docs/cli.md},
\path{letsgolegacy.secondkey/src/SecondKey.Cli/SecondKeyCli.cs},
\path{letsgolegacy.secondkey/src/SecondKey.Capture/RecordingProxy.cs},
\path{letsgolegacy.secondkey/tests/SecondKey.Cli.Tests/CaptureReplayCommandTests.cs}}

\subsection{\texttt{sk replay}}\label{sec:cli-replay}

\begin{lstlisting}
sk replay --capture .secondkey/traffic.skcap --legacy http://127.0.0.1:5080 --candidate http://127.0.0.1:5081
\end{lstlisting}

Odtwarza każdy scenariusz capture na stronie legacy, potem na kandydacie, każdą stronę przed
scenariuszem resetując do stanu początkowego --- identyczne warunki, żeby różniły się systemy,
a~nie replay. Każdy scenariusz po każdej stronie dostaje świeży słoik ciasteczek: nagrane
ciasteczka nigdy nie są wysyłane, tylko te, które strona sama ustawiła w~trakcie scenariusza.
Przekierowania nie są wykonywane; 302 to odpowiedź do porównania.

{\small
\begin{xltabular}{\linewidth}{@{}>{\raggedright\arraybackslash}p{4.9cm} >{\raggedright\arraybackslash}p{2.7cm} >{\raggedright\arraybackslash}X@{}}
\toprule
Opcja / klucz w~\texttt{secondkey.yaml} & Domyślnie & Znaczenie \\
\midrule
\endfirsthead
\toprule
Opcja / klucz w~\texttt{secondkey.yaml} & Domyślnie & Znaczenie \\
\midrule
\endhead
\bottomrule
\endlastfoot
\texttt{-{}-capture <file>}\newline \path{replay.capture} & brak --- wymagane & odtwarzany
\texttt{*.skcap} \\
\texttt{-{}-legacy <url>}\newline \path{replay.legacy.baseUrl} & brak --- wymagane & system
legacy \\
\texttt{-{}-candidate <url>}\newline \path{replay.candidate.baseUrl} & brak --- wymagane &
kandydat \\
\texttt{-{}-out <file>}\newline \path{replay.out} & \path{.secondkey/run.skrun} & przebieg \\
\texttt{-{}-scenario-mode <mode>}\newline \path{replay.scenarioMode} & \texttt{session} &
\texttt{session}: jeden scenariusz na nagraną sesję; \texttt{exchange}: jeden na wymianę \\
\texttt{-{}-timeout <seconds>}\newline \path{replay.timeoutSeconds} & 30 & jak długo czekać na
każdą odpowiedź \\
\texttt{-{}-legacy-reset-http <path>}\newline \path{replay.legacy.reset} & bez resetu & reset
strony legacy przed każdym scenariuszem przez \texttt{POST <path>}; zastępuje \texttt{reset}
z~pliku \\
\texttt{-{}-candidate-reset-http <path>}\newline \path{replay.candidate.reset} & bez resetu & to
samo dla kandydata \\
\end{xltabular}
}

\paragraph{Czego replay nie wysyła.} Z~nagranego żądania pomijane są nagłówki
\texttt{host}, \texttt{content-length}, \texttt{transfer-encoding}, \texttt{connection},
\texttt{keep-alive}, \texttt{upgrade}, \texttt{te}, \texttt{trailer},
\texttt{proxy-connection}, \texttt{proxy-authorization}, \texttt{cookie},
\texttt{accept-encoding} i~\texttt{expect}, a~także każda wartość \repopath{<redacted>}.
Ścieżka adresu bazowego strony poprzedza ścieżkę każdego żądania. Odpowiedzi są
dekompresowane, a~replay nie używa proxy systemowego.

\paragraph{Brak odpowiedzi i~podsumowanie.} Żądanie bez odpowiedzi jest zapisywane z~powodem --- \repopath{timeout}, \repopath{connection} albo \repopath{protocol} --- zamiast
odpowiedzi. Polecenie wypisuje podsumowanie:

\begin{lstlisting}
sk replay: 38 scenario(s), 580 result(s), 0 without an answer, 0 failed reset(s) -> .secondkey/run.skrun
\end{lstlisting}

i~kończy się kodem 0 albo \textbf{4}, gdy którykolwiek reset się nie powiódł. Scenariusz,
którego reset się nie udał, nie jest w~ogóle wysyłany do tej strony --- każda jego wymiana
jest zapisywana jako \repopath{reset-failed} --- więc przebieg jest zapisany, ale nie jest
uczciwym porównaniem. Brak capture albo któregoś adresu kończy się kodem 64, niepoprawny
capture albo nieustawiona zmienna środowiskowa z~connection stringiem --- kodem 3. Przebieg
przerwany przed końcem nie ma \repopath{run.end} i~komparator go nie przyjmie.

\zrodla{\path{letsgolegacy.secondkey/docs/cli.md},
\path{letsgolegacy.secondkey/src/SecondKey.Cli/Commands/ReplayCommand.cs},
\path{letsgolegacy.secondkey/src/SecondKey.Replay/ReplayEngine.cs},
\path{letsgolegacy.secondkey/docs/analysis/phase-01.md}}

\subsection{Reset stanu}\label{sec:cli-reset}

Każda strona ma w~\repopath{secondkey.yaml} jeden \repopath{reset}, wykonywany przed każdym
scenariuszem:

{\small
\begin{tabularx}{\linewidth}{@{}>{\raggedright\arraybackslash}p{3cm} >{\raggedright\arraybackslash}p{4.2cm} >{\raggedright\arraybackslash}X@{}}
\toprule
Rodzaj & Klucze & Sukces, gdy \\
\midrule
\texttt{http} & \texttt{path}; \texttt{method} (domyślnie \texttt{POST}) & strona odpowiada
kodem 2xx; \texttt{path} jest rozwiązywana względem \texttt{baseUrl} strony \\
\texttt{sqlServerSnapshot} & \texttt{connectionStringEnv}, \texttt{database},
\texttt{snapshot} & baza zostaje przywrócona ze snapshotu \\
\texttt{command} & \texttt{run}; \texttt{workingDirectory} (opcjonalny) & polecenie kończy się
kodem 0; działa pod \texttt{/bin/sh -c}, a~na Windows pod \texttt{cmd.exe /c} \\
\bottomrule
\end{tabularx}
}

Przy \repopath{http} ścieżka jest łączona z~\repopath{baseUrl} jak odnośnik względny, więc
ścieżka zaczynająca się od \repopath{/} zastępuje ścieżkę adresu bazowego. Przy
\repopath{command} szczegół porażki zawiera kod wyjścia i~wyjście błędów polecenia, a~bez
\repopath{workingDirectory} polecenie działa w~bieżącym katalogu procesu. Schemat
konfiguracji nie wymusza, że podano tylko jeden rodzaj; kod wybiera pierwszy obecny w~kolejności \repopath{http}, \repopath{sqlServerSnapshot}, \repopath{command}.

\paragraph{Ograniczenia resetu snapshotem.} \repopath{sqlServerSnapshot} przywraca bazę SQL
Server z~database snapshotu; według komentarza w~kodzie snapshoty baz danych są w~każdej
edycji SQL Server od 2016 SP1. Trzeba wiedzieć, co robi, a~czego nie:

\begin{itemize}
  \item \textbf{Stan początkowy to baza w~chwili pierwszego resetu.} Snapshot powstaje
        wtedy (pliki \texttt{<snapshot>\_<nazwa logiczna>.ss} obok plików danych) i~zostaje
        po replay; istniejący snapshot o~tej nazwie jest używany taki, jaki jest. Drugi
        replay startuje więc ze stanu, który zapisał pierwszy. Nowy stan początkowy wymaga
        usunięcia snapshotu (\repopath{DROP DATABASE [Shop\_sk]}); SQL Server odmawia też
        usunięcia bazy albo przywrócenia jej z~backupu, dopóki istnieje jej snapshot.
  \item \textbf{Każde połączenie z~bazą jest zamykane}, także pulowane połączenia aplikacji:
        przywrócenie przełącza bazę w~tryb single-user z~natychmiastowym wycofaniem
        transakcji. \repopath{sk} nie czeka, aż aplikacja dojdzie do siebie, i~nie ponawia
        pierwszego żądania scenariusza --- zapisywane jest to, co aplikacja odpowie. Pierwszą
        odpowiedź po przywróceniu trzeba raz sprawdzić, zanim się na niej polega; aplikacja,
        która uruchamia na bazie zadania cykliczne, może zawodzić do restartu.
  \item \textbf{Resetowana jest tylko baza.} To, co aplikacja trzyma w~pamięci --- cache,
        sesje, timery --- przetrwa przywrócenie. Zresetowanie tego jest zadaniem
        wywołującego; gdy nie da się tego zrobić per scenariusz, kontrakt maskuje wartości,
        na które to wpływa, i~mówi dlaczego.
  \item Login w~connection stringu musi móc tworzyć bazę i~przywracać tę bazę --- w~sandboksie \repopath{dbcreator}.
\end{itemize}

Przywrócenie ma limit 120 sekund; nieudane przywrócenie to nieudany reset z~komunikatem SQL
Server.

\zrodla{\path{letsgolegacy.secondkey/docs/cli.md},
\path{letsgolegacy.secondkey/src/SecondKey.Replay/Resets/IStateReset.cs},
\path{letsgolegacy.secondkey/src/SecondKey.Replay/Resets/SqlServerSnapshotReset.cs},
\path{letsgolegacy.secondkey/src/SecondKey.Cli/Commands/ReplayCommand.cs}}

\subsection{Sonda \texttt{sqlServer} i~korelacja}\label{sec:cli-sonda}

\paragraph{Obserwowanie bazy.} Sonda dla strony jest konfigurowana tylko w~pliku:

\begin{lstlisting}
replay:
  legacy:
    probe: { sqlServer: { connectionStringEnv: SK_LEGACY_DB, tables: [dbo.Order, dbo.ShoppingCartItem] } }
\end{lstlisting}

Sonda czyta każdy wiersz wymienionych tabel przed każdym żądaniem i~po nim i~zapisuje w~wyniku wiersze, które żądanie wstawiło, usunęło i~zmieniło (zmienione z~kluczem, stanem
przed i~po); wiersze są identyfikowane kluczem głównym, a~w~tabeli bez klucza --- całą
zawartością. Tabele bez zmian są pomijane. Sonda za każdym razem czyta całe tabele: jest
przeznaczona dla małej, zasianej bazy sandboksu, nie dla bazy wielkości produkcyjnej.
Klauzule czytają jej wynik przez korzeń \repopath{db} (sekcja~\ref{sec:formaty}).

\paragraph{Korelacja.} Wartość, którą generuje serwer, a~późniejsze żądania muszą ją odesłać
--- typowo token anti-forgery w~ukrytym polu formularza --- jest w~nagraniu nieaktualna. Reguła
korelacji (\repopath{replay.correlation}, tylko w~pliku) bierze pierwszą grupę
\repopath{regex} z~każdej odpowiedzi po tej samej stronie i~w~tym samym scenariuszu i~wpisuje
ją do pola \repopath{formField} późniejszych żądań \path{application/x-www-form-urlencoded}
albo do nagłówka \repopath{header}, jeśli nagrane żądanie ten nagłówek miało. Obowiązuje
ostatnia poznana wartość; reguła, która niczego nie poznała, niczego nie zmienia.

\zrodla{\path{letsgolegacy.secondkey/docs/cli.md},
\path{letsgolegacy.secondkey/src/SecondKey.Replay/Probes/SqlServerTableProbe.cs},
\path{letsgolegacy.secondkey/src/SecondKey.Replay/ReplayEngine.cs}}

\subsection{\texttt{sk compare}}\label{sec:cli-compare}

\begin{lstlisting}
sk compare --contract contract.yaml --run .secondkey/run.skrun --out .secondkey/verdict.json
\end{lstlisting}

{\small
\begin{tabularx}{\linewidth}{@{}>{\raggedright\arraybackslash}p{3.3cm} >{\raggedright\arraybackslash}p{3.3cm} >{\raggedright\arraybackslash}X@{}}
\toprule
Opcja & W~\texttt{secondkey.yaml} & Domyślnie \\
\midrule
\texttt{-{}-contract <file>} & \texttt{compare.contract} & brak --- wymagane \\
\texttt{-{}-run <file>} & \texttt{compare.run} & \path{.secondkey/run.skrun} \\
\texttt{-{}-out <file>} & \texttt{compare.out} & \path{.secondkey/verdict.json} \\
\bottomrule
\end{tabularx}
}

Umieszcza każdą odtworzoną wymianę w~jednej z~czterech klas --- \repopath{equal},
\repopath{equal-under-contract}, \repopath{regression}, \repopath{fix-candidate} --- i~zapisuje werdykt (klasy, kolejność i~format: sekcja~\ref{sec:formaty}). Kontrakt i~przebieg
są walidowane w~całości; werdykt dostaje czas z~dokładnością do sekundy i~nazwy plików
wejściowych bez katalogów, a~przed zapisem jest sprawdzany własnym schematem. Polecenie
wypisuje liczby i~każdą wymianę, na którą musi spojrzeć osoba:

\begin{itemize}
  \item \texttt{sk compare: <outcome> — <n> exchange(s): <e> equal, <u> equal under contract,
        <r> regression(s), <f> fix candidate(s) -> <plik>};
  \item po jednej linii na regresję i~kandydata na poprawkę: klasa, identyfikator wymiany,
        metoda ze ścieżką i~zapytaniem, powody;
  \item \texttt{clauses: <total> (<a> accepted, <x> exercised, <y> unexercised)} z~dopiskiem
        \texttt{absence share <p>\%}, gdy jest;
  \item po jednej linii na zaakceptowaną klauzulę o~statusie \repopath{regressed},
        \repopath{fixed} albo \repopath{unexercised}.
\end{itemize}

Kody wyjścia: \textbf{0} przy \repopath{pass}, \textbf{1} przy \repopath{fail} (jakakolwiek
regresja), \textbf{2} przy \repopath{review} (bez regresji, kandydat na poprawkę do decyzji).
Przebieg bez wyników jest niepoprawnym wejściem (\textbf{3}): nic nie porównano, więc nic nie
może przejść; kod 3 dają też niepoprawny kontrakt albo przebieg. Brak kontraktu to kod 64.

\zrodla{\path{letsgolegacy.secondkey/docs/cli.md},
\path{letsgolegacy.secondkey/src/SecondKey.Cli/Commands/CompareCommand.cs},
\path{letsgolegacy.secondkey/tests/SecondKey.Cli.Tests/CompareCommandTests.cs}}

\subsection{\texttt{sk gate}}\label{sec:cli-gate}

\begin{lstlisting}
sk gate --sarif portcullis.sarif
\end{lstlisting}

\repopath{-{}-sarif} przyjmuje jeden albo więcej logów SARIF 2.1.0 (np. z~Portcullis); w~pliku --- lista \repopath{gate.sarif}. Polecenie podsumowuje logi per reguła i~kończy się kodem
\textbf{1}, gdy znalezisko blokuje, a~0 w~przeciwnym razie. \finding{Błąd blokuje, chyba że
obejmuje go zaakceptowane wyciszenie albo log oznacza go jako znany z~baseline
(\texttt{baselineState} \texttt{unchanged} albo \texttt{absent}).} Ostrzeżenia i~notatki są
raportowane i~nigdy nie blokują. Wyciszenie jest zaakceptowane, gdy nie ma statusu albo ma
status \repopath{accepted}. Poziom wyniku pochodzi z~samego wyniku, potem z~domyślnego poziomu
reguły, a~w~ostateczności jest ostrzeżeniem; wynik o~rodzaju (\repopath{kind}) innym niż
\repopath{fail} nie ma poziomu. Reguły są zbierane z~drivera i~rozszerzeń narzędzia; raport w~pakiecie wymienia każdą zadeklarowaną regułę, także te bez znalezisk --- zapisuje, co
sprawdzono, a~nie tylko, co znaleziono --- a~wyjście polecenia wymienia reguły, które coś
znalazły. Plik, który nie jest SARIF 2.1.0, jest niepoprawnym wejściem (kod 3), nigdy nie
jest czytany jako „brak znalezisk”. Bez żadnego logu polecenie kończy się kodem 64.

Wyjście: \texttt{sk gate: pass|fail — <b> blocking of <n> finding(s) in <k> log(s)}, potem
wiersz na regułę z~liczbą błędów, ostrzeżeń, notatek, wyciszonych i~blokujących, a~na końcu
linia \texttt{blocks:} z~położeniem, regułą i~komunikatem każdego blokującego znaleziska.
Syntetyczny log \path{samples/gate/portcullis.sarif} ma dwa ostrzeżenia i~wyciszony błąd, więc
bramka przechodzi.

\zrodla{\path{letsgolegacy.secondkey/docs/formats/evidence.md},
\path{letsgolegacy.secondkey/docs/cli.md},
\path{letsgolegacy.secondkey/src/SecondKey.Cli/Commands/GateCommand.cs},
\path{letsgolegacy.secondkey/src/SecondKey.Evidence/Sarif/SarifLog.cs},
\path{letsgolegacy.secondkey/src/SecondKey.Evidence/Sarif/GateSummary.cs},
\path{letsgolegacy.secondkey/samples/gate/portcullis.sarif}}

\subsection{\texttt{sk evidence}}\label{sec:cli-evidence}

\begin{lstlisting}
sk evidence --verdict verdict.json --contract contract.yaml --run run.skrun \
            --sarif portcullis.sarif --out evidence --pdf auto
\end{lstlisting}

{\small
\begin{tabularx}{\linewidth}{@{}>{\raggedright\arraybackslash}p{3.3cm} >{\raggedright\arraybackslash}p{3.8cm} >{\raggedright\arraybackslash}X@{}}
\toprule
Opcja & W~\texttt{secondkey.yaml} & Domyślnie i~znaczenie \\
\midrule
\texttt{-{}-verdict <file>} & \texttt{evidence.verdict} & \path{.secondkey/verdict.json} \\
\texttt{-{}-contract <file>} & \texttt{evidence.contract}, potem \texttt{compare.contract} &
wymagane: kontrakt, z~którego policzono werdykt \\
\texttt{-{}-run <file>} & \texttt{evidence.run} & gdy nie podano, przebieg nie trafia do
pakietu; podany pozwala każdemu przeliczyć werdykt \\
\texttt{-{}-sarif <file>...} & \texttt{evidence.sarif} & logi bramki; bez nich pakiet mówi, że
nie dostarczono wyników analizy statycznej \\
\texttt{-{}-out <dir>} & \texttt{evidence.out} & \path{.secondkey/evidence} \\
\texttt{-{}-pdf <mode>} & \texttt{evidence.pdf} & \texttt{auto}; także \texttt{required},
\texttt{off} \\
\bottomrule
\end{tabularx}
}

Polecenie kończy się kodem 0, gdy pakiet został zapisany, \emph{niezależnie od werdyktu}:
kod wyjścia werdyktu należy do \repopath{sk compare}, a~pipeline powinien opublikować pakiet,
który wyjaśnia porażkę. Odrzuca kodem 3 kontrakt albo przebieg, który nie jest tym, który
nazywa werdykt, a~także niepoprawny werdykt albo log SARIF. Katalog wyjściowy z~poprzednim
pakietem jest nadpisywany dokładnie w~zakresie plików z~jego manifestu; katalog z~czymkolwiek
innym jest odrzucany (kod 64, „… holds no evidence pack; choose an empty or new directory”).
Wyjście wymienia katalog, pliki pakietu oraz linię \texttt{behaviour: <outcome> · gate:
<wynik> · PDF: <stan>}.

\paragraph{PDF.} \repopath{-{}-pdf auto} (domyślnie) drukuje raport, gdy na maszynie jest
przeglądarka oparta na Chromium, a~gdy jej nie ma, mówi o~tym i~zapisuje pakiet bez PDF;
\repopath{-{}-pdf required} czyni brak PDF porażką (kod 4); \repopath{-{}-pdf off} pomija PDF.
Wartość spoza tych trzech --- w~opcji albo w~\repopath{evidence.pdf} --- kończy się kodem 64.
Przeglądarka jest szukana kolejno:

\begin{enumerate}
  \item \repopath{SK\_CHROME\_PATH} --- jeśli jest ustawiona, plik musi istnieć, inaczej to
        błąd, a~nie ciche przejście dalej;
  \item \texttt{PATH}, katalog po katalogu: \texttt{google-chrome},
        \texttt{google-chrome-stable}, \texttt{chrome}, \texttt{microsoft-edge},
        \texttt{microsoft-edge-stable}, \texttt{msedge}, \texttt{chromium},
        \texttt{chromium-browser} (na Windows z~\texttt{.exe}) --- Google Chrome i~Edge przed
        Chromium, którego buildy bardziej się różnią;
  \item miejsca instalacji: na Windows Chrome i~Edge w~\texttt{ProgramFiles},
        \texttt{ProgramFiles(x86)} i~\texttt{LOCALAPPDATA}; na macOS aplikacje Google Chrome,
        Microsoft Edge i~Chromium; przeglądarki Playwright z~\repopath{PLAYWRIGHT\_BROWSERS\_PATH} albo \path{/opt/pw-browsers}, od najnowszej.
\end{enumerate}

Przeglądarka, która nie wydrukowała w~ciągu minuty, jest zatrzymywana, a~błąd podaje ostatnią
rzecz, którą napisała. Działa headless, z~własnym tymczasowym profilem i~z~wyłączonym
sandboksem, bo sandbox jest niedostępny dla roota w~kontenerach i~na niektórych obrazach CI;
renderuje wyłącznie własną statyczną stronę pakietu, która nie może uruchomić skryptu ani
wykonać żądania. Plik, który nie zaczyna się od \texttt{\%PDF-}, nie liczy się jako PDF.

\zrodla{\path{letsgolegacy.secondkey/docs/formats/evidence.md},
\path{letsgolegacy.secondkey/docs/cli.md},
\path{letsgolegacy.secondkey/src/SecondKey.Cli/Commands/EvidenceCommand.cs},
\path{letsgolegacy.secondkey/src/SecondKey.Evidence/EvidencePack.cs},
\path{letsgolegacy.secondkey/src/SecondKey.Evidence/PdfRenderer.cs},
\path{letsgolegacy.secondkey/tests/SecondKey.Cli.Tests/GateEvidenceCommandTests.cs}}

\subsection{\texttt{sk validate}, \texttt{sk mine} i~\texttt{sk mutate}}\label{sec:cli-validate}

\repopath{sk validate <file>...} sprawdza jeden albo więcej plików względem ich formatów i~wypisuje każdy błąd z~położeniem. Rodzaj pliku rozpoznaje po nazwie: \repopath{secondkey.yaml}
jest walidowany jako konfiguracja, \repopath{*.skcap} i~\repopath{*.skrun} --- jako
capture i~przebieg, \repopath{*.yaml} i~\repopath{*.yml} --- jako kontrakt, a~\repopath{*.json} --- jako werdykt, gdy jego \repopath{kind} to
\repopath{secondkey.verdict} (albo gdy nie jest poprawnym JSON-em). Każdy inny plik dostaje
błąd „not a~Second Key artifact: expected *.skcap, *.skrun, a~contract *.yaml or a~verdict
*.json”, a~nieistniejący --- „no such file”. Dla każdego pliku wypisywane jest
\repopath{<plik>: valid} albo liczba błędów i~ich lista. Kod 0, gdy wszystkie pliki są
poprawne, w~przeciwnym razie 3.

\repopath{sk mine} (C2) i~\repopath{sk mutate} (C9) istnieją w~planie, ale nie w~tej wersji.
Czytają konfigurację, wypisują na przykład „sk mine: not available in this build --- planned
under C2, phase 02 (docs/WORKPLAN.md).” i~kończą się kodem \textbf{70}.

\zrodla{\path{letsgolegacy.secondkey/src/SecondKey.Cli/Commands/ValidateCommand.cs},
\path{letsgolegacy.secondkey/src/SecondKey.Cli/Commands/StubCommand.cs},
\path{letsgolegacy.secondkey/src/SecondKey.Cli/SecondKeyCli.cs},
\path{letsgolegacy.secondkey/docs/formats/README.md}}

\subsection{Kody wyjścia}\label{sec:cli-kody}

Pipeline rozgałęzia się na tych kodach; nie zmieniają się bez zmiany wersji głównej, a~klasa
\repopath{ExitCodes} nazywa zmianę któregokolwiek zmianą łamiącą każdy pipeline, który go
używa. Test sprawdza, że kody są różne i~równe dokładnie 0, 1, 2, 3, 4, 64 i~70.

{\small
\begin{tabularx}{\linewidth}{@{}l l >{\raggedright\arraybackslash}X@{}}
\toprule
Kod & Nazwa & Znaczenie \\
\midrule
0 & \texttt{Success} & werdykt przeszedł, bramka przeszła, artefakt jest poprawny \\
1 & \texttt{Failed} & to, co sprawdzano, nie przeszło: regresja w~werdykcie, blokujące
znalezisko w~bramce \\
2 & \texttt{Review} & brak regresji, ale kandydat na poprawkę wymaga decyzji osoby \\
3 & \texttt{InvalidInput} & wejście --- artefakt albo konfiguracja --- jest niepoprawne;
niczego nie rozstrzygnięto \\
4 & \texttt{RuntimeError} & polecenie nie mogło się zakończyć: system nieosiągalny, plik nie
do zapisania \\
64 & \texttt{Usage} & sam wiersz poleceń jest błędny \\
70 & \texttt{NotImplemented} & polecenie istnieje w~planie, ale nie w~tej wersji \\
\bottomrule
\end{tabularx}
}

Polecenie zatrzymane Ctrl+C albo SIGTERM ma minutę na czyste zakończenie; jeśli się nie
zatrzyma, jest kończone kodem sygnału, 130 albo 143. W~kodzie poszczególne rodzaje porażek
mapują się tak: błąd walidacji artefaktu i~błąd konfiguracji (brak jawnie wskazanego pliku,
niepoprawny YAML, nieznany klucz, nieobsługiwana wersja, nieustawiona zmienna z~sekretem) →
3; błąd wiersza poleceń, brak wymaganego wejścia, adres, który nie jest bezwzględnym URL
http(s) --- także w~\repopath{secondkey.yaml} --- i~wartość \repopath{evidence.pdf} spoza
dozwolonych → 64; anulowanie → 4 z~komunikatem „cancelled”; każdy inny wyjątek, np. brak
pliku wejściowego, → 4 z~komunikatem \texttt{error: …}.

\zrodla{\path{letsgolegacy.secondkey/docs/cli.md},
\path{letsgolegacy.secondkey/src/SecondKey.Cli/ExitCodes.cs},
\path{letsgolegacy.secondkey/src/SecondKey.Cli/SecondKeyCli.cs},
\path{letsgolegacy.secondkey/tests/SecondKey.Cli.Tests/CommandLineTests.cs},
\path{letsgolegacy.secondkey/tests/SecondKey.Cli.Tests/GateEvidenceCommandTests.cs}}

\subsection{Plik \texttt{secondkey.yaml}}\label{sec:cli-konfiguracja}

\repopath{secondkey.yaml} daje wartości domyślne każdemu poleceniu, żeby krok pipeline'u był
wywołaniem \repopath{sk replay}, a~nie linią opcji. Przykład z~\path{docs/cli.md}:

\begin{lstlisting}
version: 1
capture:
  listen: http://127.0.0.1:8080          # where the recording proxy listens
  target: http://127.0.0.1:5080          # the legacy system behind it
  out: .secondkey/traffic.skcap
  sessionKeys: [ASP.NET_SessionId]       # cookies that identify a user session
  redactHeaders: [authorization, cookie, set-cookie]
replay:
  capture: .secondkey/traffic.skcap
  out: .secondkey/run.skrun
  scenarioMode: session                  # session | exchange
  legacy:
    baseUrl: http://127.0.0.1:5080
    reset: { sqlServerSnapshot: { connectionStringEnv: SK_LEGACY_DB, database: Shop, snapshot: Shop_sk } }
  candidate:
    baseUrl: http://127.0.0.1:5081
    reset: { http: { path: /__sk/reset } }
  correlation:                           # values the server generates that later requests must echo
    - { name: csrf, regex: 'name="__RequestVerificationToken" type="hidden" value="([^"]+)"', formField: __RequestVerificationToken }
compare: { contract: contract.yaml, run: .secondkey/run.skrun, out: .secondkey/verdict.json }
gate: { sarif: [gate.sarif] }
evidence: { verdict: .secondkey/verdict.json, contract: contract.yaml, run: .secondkey/run.skrun, sarif: [gate.sarif], out: .secondkey/evidence, pdf: auto }
\end{lstlisting}

{\small
\begin{xltabular}{\linewidth}{@{}>{\raggedright\arraybackslash}p{5.2cm} >{\raggedright\arraybackslash}X@{}}
\toprule
Klucz & Znaczenie \\
\midrule
\endfirsthead
\toprule
Klucz & Znaczenie \\
\midrule
\endhead
\bottomrule
\endlastfoot
\texttt{version} & wersja pliku; domyślnie 1, inna wartość to błąd („version 2 is not
supported; this build reads version 1”) \\
\texttt{capture.listen}, \texttt{target}, \texttt{out}, \texttt{sessionKeys},
\texttt{redactHeaders}, \texttt{maxBodyBytes} & opcje \texttt{sk capture}
(sekcja~\ref{sec:cli-capture}) \\
\texttt{replay.capture}, \texttt{out}, \texttt{scenarioMode}, \texttt{timeoutSeconds} & opcje
\texttt{sk replay} (sekcja~\ref{sec:cli-replay}) \\
\texttt{replay.legacy}, \texttt{replay.candidate} & każda strona: \texttt{baseUrl},
\texttt{reset}, \texttt{probe} \\
\texttt{…reset.http} & \texttt{method} (domyślnie \texttt{POST}), \texttt{path} (wymagany) \\
\texttt{…reset.sqlServerSnapshot} & \texttt{connectionStringEnv}, \texttt{database},
\texttt{snapshot} (wszystkie wymagane) \\
\texttt{…reset.command} & \texttt{run} (wymagany), \texttt{workingDirectory} \\
\texttt{…probe.sqlServer} & \texttt{connectionStringEnv}, \texttt{tables} (oba wymagane) \\
\texttt{replay.correlation[]} & \texttt{name}, \texttt{regex} (wymagane; pierwsza grupa to
wartość), \texttt{formField}, \texttt{header} \\
\texttt{compare.contract}, \texttt{run}, \texttt{out} & opcje \texttt{sk compare} \\
\texttt{gate.sarif} & lista logów SARIF dla \texttt{sk gate} \\
\texttt{evidence.verdict}, \texttt{contract}, \texttt{run}, \texttt{sarif}, \texttt{out},
\texttt{pdf} & opcje \texttt{sk evidence} \\
\end{xltabular}
}

\paragraph{Ścisłość.} Plik jest ścisły: nieznany klucz jest błędem (np. literówka
\repopath{resett} albo \repopath{scenarioModes}), a~brak wymaganego klucza w~sekcji resetu
też. Pusty plik to pusta konfiguracja. YAML jest typowany tak samo jak w~kontrakcie (YAML 1.2
core schema). \repopath{sk validate secondkey.yaml} sprawdza plik bez uruchamiania czegokolwiek.
Wartość \repopath{scenarioMode} nie jest w~pliku sprawdzana: każda inna niż
\repopath{exchange} oznacza \repopath{session}; opcja \repopath{-{}-scenario-mode} przyjmuje
tylko te dwie wartości.

\paragraph{Sekrety.} \finding{Sekrety nigdy nie mieszkają w~tym pliku.} Connection string
jest wskazywany nazwą zmiennej środowiskowej, która go przechowuje
(\repopath{connectionStringEnv}), a~\path{secrets.env.example} wymienia każdą taką zmienną.
Zmienna nieustawiona albo pusta zatrzymuje polecenie z~jej nazwą („environment variable
SK\_LEGACY\_DB is not set (see secrets.env.example)”) i~kodem 3.

\paragraph{Ścieżki.} Ścieżki względne rozwiązuje się względem katalogu, w~którym leży plik.
W~kodzie wyjątkiem jest \repopath{reset.command.workingDirectory}, przekazywany procesowi bez
rozwiązywania. Na Windows ścieżkę pisze się z~ukośnikami (\path{C:/sk/traffic.skcap}) albo w~apostrofach (\path{'C:\sk\traffic.skcap'}): w~cudzysłowie YAML czyta ukośnik wsteczny jako
początek sekwencji ucieczki, więc \path{"C:\sk\a.skcap"} nie jest ścieżką, na którą wygląda.

\zrodla{\path{letsgolegacy.secondkey/docs/cli.md},
\path{letsgolegacy.secondkey/src/SecondKey.Cli/Configuration/SecondKeyConfig.cs},
\path{letsgolegacy.secondkey/src/SecondKey.Cli/Commands/ReplayCommand.cs},
\path{letsgolegacy.secondkey/tests/SecondKey.Cli.Tests/ConfigurationTests.cs},
\path{letsgolegacy.secondkey/secrets.env.example}}

\subsection{Demonstracja end-to-end na sample shop}\label{sec:cli-e2e}

Dwa polecenia z~README uruchamiają cały łańcuch lokalnie:

\begin{lstlisting}
./scripts/setup.sh     # once: checks the tools, installs the secret-scanning commit hook
./scripts/e2e.sh       # the whole chain on the sample shop
\end{lstlisting}

\paragraph{Sample shop.} \path{samples/SampleShop} to mała aplikacja ASP.NET Core, która gra
obie role demonstracji fazy 01. Wariant wybiera \repopath{SAMPLESHOP\_VARIANT}
(\repopath{legacy} domyślnie albo \repopath{candidate}); sesję niesie ciasteczko
\repopath{sk\_session}. Punkty końcowe: \texttt{GET /products?sort=name} (JSON),
\texttt{POST /cart/items}, \texttt{POST /cart/coupon} i~\texttt{POST /checkout} (JSON),
\texttt{GET /cart} (HTML), \texttt{GET /health} (zwraca nazwę wariantu) oraz
\texttt{POST /\_\_sk/reset}, który czyści koszyki i~odpowiada 204 tylko przy
\repopath{SAMPLESHOP\_ALLOW\_RESET=1}, a~w~przeciwnym razie 404. Kandydat różni się od strony
legacy dokładnie tak, jak prawdziwa migracja, więc każda klasa werdyktu ma co znaleźć:

\begin{itemize}
  \item lista produktów jest sortowana bieżącą kulturą; strona legacy działa z~\path{DOTNET_SYSTEM_GLOBALIZATION_INVARIANT=1} i~kolejność się zmienia, choć kod
        nie --- pułapka NLS → ICU migracji z~.NET Framework, odtworzona na Linuksie;
  \item sumy wierszy u~kandydata pochodzą z~arytmetyki \texttt{double} (szum, który pochłania
        tolerancja);
  \item kupon większy niż koszyk sprowadza sumę kandydata poniżej zera (regresja);
  \item zamówienie z~pustym koszykiem wywraca stronę legacy (500), a~kandydat odmawia czysto
        (400) --- kandydat na poprawkę;
  \item identyfikatory zamówień i~znaczniki czasu są nowe przy każdym przebiegu (pochłaniają
        je maski); strona HTML koszyka ma u~obu wariantów inny markup, który kontraktu nie
        obchodzi.
\end{itemize}

\paragraph{Przebieg \texttt{scripts/e2e.sh}.} Skrypt buduje solution w~konfiguracji Release,
uruchamia stronę legacy (niezmienna globalizacja) i~kandydata (ICU), czeka na
\texttt{/health} obu, uruchamia \repopath{sk capture} przed stroną legacy z~\repopath{-{}-session-key sk\_session} i~\repopath{-{}-exit-after 10}, wysyła przez proxy
dziesięć żądań, czeka, aż capture zatrzyma się sam, i~waliduje nagranie:

\begin{lstlisting}
call GET  "/products?sort=name" none                                          # a visitor looks at the catalogue
call POST /cart/items           buyer   '{"sku":"ECL-02","quantity":3}'       # a buyer orders
call GET  /cart                 buyer
call POST /checkout             buyer   '{"email":"buyer@example.test"}'
call POST /cart/items           ruler   '{"sku":"ANG-04","quantity":9}'       # 9 x 3.10 is 27.900000000000002 in double arithmetic
call GET  /cart                 ruler
call POST /cart/items           bargain '{"sku":"APL-01","quantity":1}'       # a coupon larger than the cart
call POST /cart/coupon          bargain '{"code":"BIGSALE"}'
call GET  /cart                 bargain
call POST /checkout             empty   '{}'                                  # checkout with nothing in the cart
\end{lstlisting}

Trzecia kolumna to słoik ciasteczek, czyli sesja użytkownika; \texttt{none} to gość bez
sesji. Potem skrypt odtwarza capture na obu stronach z~resetem
\texttt{-{}-legacy-reset-http /\_\_sk/reset} i~\texttt{-{}-candidate-reset-http
/\_\_sk/reset}, uruchamia \repopath{sk compare} z~\path{samples/contract.yaml} i~wymaga kodu 1
(próbka musi oblać), uruchamia \repopath{sk gate} na \path{samples/gate/portcullis.sarif}
(ostrzeżenia i~wyciszony błąd, nic blokującego --- kod 0) i~\repopath{sk evidence}. Wszystko
trafia do \repopath{\$E2E\_DIR}, domyślnie \path{.secondkey/e2e}; porty ustawiają
\repopath{E2E\_LEGACY\_PORT} (5080), \repopath{E2E\_CANDIDATE\_PORT} (5081) i~\repopath{E2E\_PROXY\_PORT} (8080), a~tryb PDF --- \repopath{SK\_PDF} (domyślnie
\repopath{auto}; CI ustawia \repopath{required}). Kontrakt demonstracji celowo nie ma żadnej
reguły o~kolejności listy produktów: kandydat sortuje te same nazwy inaczej, nic nie mówi,
że kolejność nie ma znaczenia, więc różnica musi wyjść jako regresja, której nikt nie
wyjaśnił.

\paragraph{Co skrypt sprawdza.} Na końcu skrypt w~Pythonie sprawdza, że werdykt zawiera to,
do czego próbka jest zbudowana:

\begin{checklista}
  \item każda z~czterech klas występuje co najmniej raz, a~wynik to \repopath{fail};
  \item jedyna wymiana \texttt{/products} jest regresją z~powodem wyłącznie
        \repopath{uncovered-difference}, a~różnica w~\path{response.body.json.items[0]} nie
        ma \repopath{normalizedBy} (zmiana kolejności NLS → ICU);
  \item co najmniej dwie wymiany mają \repopath{clause-regression} dla
        \repopath{TOTAL-NEVER-NEGATIVE} albo \repopath{CART-PAGE-TOTAL-NEVER-NEGATIVE}
        (skumulowany rabat łamie obie klauzule o~sumie);
  \item kandydat na poprawkę ma powód dokładnie \texttt{clause-fixed: CHECKOUT-NEVER-5XX}
        (pusty koszyk);
  \item wymiany równe pod kontraktem mają powody \texttt{normalized: masks.guids} i~\texttt{normalized: masks.timestamps} oraz jakiś
        \texttt{normalized: tolerances: …} (arytmetyka \texttt{double});
  \item skrót każdego pliku zgadza się z~manifestem, a~pakiet zawiera \repopath{index.html},
        \repopath{verdict.json}, \repopath{contract.yaml}, \repopath{run.skrun},
        \path{statement.intoto.json}, \path{sarif/portcullis.sarif} i~--- przy
        \repopath{SK\_PDF=required} --- \repopath{report.pdf}.
\end{checklista}

Job \repopath{e2e} w~CI uruchamia ten sam skrypt i~publikuje pakiet jako artefakt
\repopath{evidence-pack} (sekcja~\ref{sec:rdzen}). README podsumowuje, co próbka pokazuje:
to, co ukrywa prawdziwa migracja --- listę produktów posortowaną inaczej, na co nic w~kontrakcie nie pozwala; rabat, który sprowadza sumę poniżej zera; i~zamówienie, które przestało
się wywracać na pustym koszyku, przedstawione osobie jako poprawka.

\zrodla{\path{letsgolegacy.secondkey/README.md},
\path{letsgolegacy.secondkey/scripts/e2e.sh},
\path{letsgolegacy.secondkey/samples/SampleShop/Program.cs},
\path{letsgolegacy.secondkey/samples/contract.yaml},
\path{letsgolegacy.secondkey/samples/gate/portcullis.sarif},
\path{letsgolegacy.secondkey/.github/workflows/e2e.yml}}

\subsection{Jak czytać wynik}\label{sec:cli-wynik}

\begin{enumerate}
  \item \textbf{Kod wyjścia \texttt{sk compare}} rozstrzyga krok pipeline'u: 0 --- zachowanie
        dotrzymuje kontraktu, 1 --- jest regresja, 2 --- jest kandydat na poprawkę do
        decyzji. \repopath{sk gate} daje osobno 0 albo 1 dla kodu, a~\repopath{sk evidence}
        zawsze 0, gdy pakiet powstał.
  \item \textbf{Wyjście \texttt{sk compare}} wymienia tylko to, na co musi spojrzeć osoba:
        regresje, kandydatów na poprawkę i~zaakceptowane klauzule \repopath{regressed},
        \repopath{fixed} albo \repopath{unexercised}. Skrypt e2e wypisuje dodatkowo każdą
        wymianę (klasa, identyfikator, żądanie, powody), podsumowanie w~\repopath{\$GITHUB\_STEP\_SUMMARY} i~na końcu \texttt{e2e: ok — <n> exchanges, every
        class present, pack of <m> files at <katalog>}.
  \item \textbf{Pakiet dowodowy} otwiera się od \repopath{index.html}: odznaki „Behaviour”
        i~„Gate”, kafelki podsumowania, potem regresje i~kandydaci na poprawkę z~różnicami i~klauzulami, które złamała któraś strona, tabela klauzul, to, co kontrakt uznaje za
        nieistotne, i~to, czego dowód nie pokazuje --- ruch, którego nagranie nie
        zawierało, zaakceptowane klauzule \repopath{unexercised}, klauzule
        \repopath{proposed}, klauzule \repopath{violated-both} i~brak podpisu.
  \item \textbf{Każda regresja} kończy się przeglądaną zmianą: poprawką kandydata albo zmianą
        kontraktu --- regułą normalizacji, jeśli różnica nie ma znaczenia, albo klauzulą, jeśli
        ma. \textbf{Każdy kandydat na poprawkę} wymaga decyzji osoby, czy zmiana jest
        pożądana. Klauzula \repopath{violated-both} to defekt, który migracja zachowała; nie
        oblewa przebiegu.
  \item \textbf{Integralność} sprawdza się niezależnie od Second Key: skróty z~\repopath{manifest.json}, walidacja \repopath{verdict.json} i~kontraktu oraz ponowne
        \repopath{sk compare} na \repopath{run.skrun} z~pakietu, które daje ten sam werdykt,
        poza \repopath{createdAt} (sekcja~\ref{sec:formaty}).
\end{enumerate}

\zrodla{\path{letsgolegacy.secondkey/docs/cli.md},
\path{letsgolegacy.secondkey/docs/formats/verdict.md},
\path{letsgolegacy.secondkey/docs/formats/evidence.md},
\path{letsgolegacy.secondkey/src/SecondKey.Cli/Commands/CompareCommand.cs},
\path{letsgolegacy.secondkey/src/SecondKey.Evidence/HtmlReport.cs},
\path{letsgolegacy.secondkey/scripts/e2e.sh}}
